% =====================================================================
\section{Le socle transverse}

Cinq mécanismes servent à toutes les fonctionnalités. Les connaître évite de
répéter la même explication douze fois pendant la soutenance.

\subsection{Authentification et rôles}

L'authentification s'appuie sur la table \texttt{SECURITY\_USERS}, portée par
l'entité \texttt{SecurityUser} et interrogée par \texttt{SecurityService}. Ce
service est le seul \texttt{@Singleton} du projet, et il est marqué
\texttt{@Startup} : une instance unique, créée au démarrage du serveur.

\begin{lstlisting}[style=java]
@Singleton
@Startup
public class SecurityService {
    @PersistenceContext
    private EntityManager em;

    public SecurityUser authenticate(String username, String password) {
        TypedQuery<SecurityUser> query = em.createQuery(
            "SELECT s FROM SecurityUser s WHERE s.username = :username "
          + "AND s.password = :password", SecurityUser.class);
        query.setParameter("username", username);
        query.setParameter("password", password);
        List<SecurityUser> results = query.getResultList();
        return results.isEmpty() ? null : results.get(0);
    }
}
\end{lstlisting}

Le \texttt{LoginServlet} appelle ce service, puis, en cas de succès, place
trois attributs en session --- et rien d'autre :

\begin{lstlisting}[style=java]
HttpSession session = request.getSession();
session.setAttribute("user",   username);
session.setAttribute("role",   role);
session.setAttribute("userId", businessUser.getId());
response.sendRedirect(request.getContextPath() + "/home");
\end{lstlisting}

\textbf{Trois rôles} circulent dans l'application, exactement ceux stockés dans
\texttt{SECURITY\_USERS} :

\noindent\begin{tabularx}{\textwidth}{|>{\columncolor{btkpale}\ttfamily}L{4.2cm}|X|}
\hline
\textmd{\textbf{Rôle}} & \textbf{Périmètre}\\
\hline
ADMIN & Tout le périmètre fonctionnel\\
\hline
DIRECTEUR\_COMMERCIAL & Soumet les demandes de création, valide les demandes de modification, consulte\\
\hline
USER & Son pointage, ses demandes, ses notifications\\
\hline
\end{tabularx}

\vspace{0.6em}
Le contrôle est fait \textbf{dans chaque servlet}, au début de \texttt{doGet}
et de \texttt{doPost}, selon un motif constant :

\begin{lstlisting}[style=java]
HttpSession session = request.getSession(false);
if (session == null
        || session.getAttribute("user") == null
        || !"ADMIN".equals(session.getAttribute("role"))) {
    response.sendRedirect(request.getContextPath() + "/index.jsp");
    return;
}
\end{lstlisting}

\begin{retenir}
\textbf{Pourquoi \texttt{getSession(false)} ?} Avec \texttt{true} (le défaut),
l'appel \emph{crée} une session si elle n'existe pas --- le test
\enquote{la session est-elle nulle} ne détecterait alors plus jamais un
visiteur non connecté. Avec \texttt{false}, la méthode renvoie \texttt{null}
quand aucune session n'existe : c'est ce qui permet le contrôle.
\end{retenir}

\subsection{Le filtre de notifications}

\texttt{NotificationFilter} est annoté \texttt{@WebFilter("/*")} : il s'exécute
\textbf{avant tout servlet}, sur toutes les URL. Son rôle est d'alimenter les
compteurs affichés dans le bandeau, sans que chaque servlet ait à le refaire.

\begin{lstlisting}[style=java]
if (session != null) {
    String role = (String) session.getAttribute("role");
    if ("ADMIN".equals(role)) {
        request.setAttribute("pendingClientDemandes",  demandeClientService.countPending());
        request.setAttribute("pendingEmployeDemandes", demandeEmployeService.countEnAttente());
        request.setAttribute("approvedModifs",         demandeModificationService.countInCours());
    } else if ("DIRECTEUR_COMMERCIAL".equals(role)) {
        request.setAttribute("pendingModifs", demandeModificationService.countPending());
    }
}
chain.doFilter(request, response);
\end{lstlisting}

Les compteurs dépendent du rôle : l'administrateur voit ce qu'il doit
\emph{exécuter}, le directeur ce qu'il doit \emph{valider}.

\subsection{Le journal d'audit}

\texttt{JournalService} écrit dans \texttt{JOURNAL\_ACTION} (entité
\texttt{JournalAction}). Chaque écriture sensible appelle \texttt{log(\dots)}
avec l'auteur, son rôle, l'action, le type et l'identifiant de l'objet touché,
et un libellé. Exemple relevé dans \texttt{DemandeAgenceAdminServlet} :

\begin{lstlisting}[style=java]
journalService.log(userId, role,
        "APPROVE_DEMANDE_AGENCE", "DEMANDE_AGENCE", demande.getId(),
        "Approbation de la demande d'agence code=" + demande.getCodeAgencePropose());
\end{lstlisting}

\subsection{Les notifications}

\texttt{NotificationService} écrit dans \texttt{NOTIFICATION}. Une notification
vise \textbf{un utilisateur} et porte un type et un message :

\begin{lstlisting}[style=java]
notificationService.notifyUser(demande.getDirecteur(),
        "DEMANDE_AGENCE_VALIDEE",
        "Votre demande d'agence " + demande.getCodeAgencePropose() + " a ete approuvee.");
\end{lstlisting}

C'est le mécanisme qui referme chaque circuit de demande : celui qui a soumis
est prévenu de la décision.

\subsection{Le gabarit de page}

Deux fragments sont inclus par toutes les pages authentifiées :

\begin{itemize}[leftmargin=*,itemsep=1pt]
  \item \texttt{WEB-INF/sidebar.jsp} --- le menu latéral. Les entrées sont
        conditionnées par le rôle avec \texttt{<c:if>} ; le paramètre
        \texttt{activePage} met en relief l'entrée courante.
  \item \texttt{WEB-INF/top-bar.jsp} --- le bandeau, qui affiche les compteurs
        posés par le filtre et oriente le clic selon le rôle :
\end{itemize}

\begin{lstlisting}[style=jsp]
<c:set var="notifLink"
       value="${sessionScope.role == 'ADMIN' ? 'demandes-admin' : 'gestion-demandes'}" />
\end{lstlisting}

\subsection{Les seize entités et leurs tables}

\noindent\begin{tabularx}{\textwidth}{|>{\ttfamily}L{4.3cm}|>{\ttfamily}L{4.3cm}|X|}
\hline
\textmd{\textbf{Entité}} & \textmd{\textbf{Table Oracle}} & \textbf{Ce qu'elle porte}\\
\hline
SecurityUser & SECURITY\_USERS & comptes et rôles\\
\hline
Utilisateur & B\_UTILISATEURS & employés (dont les gestionnaires)\\
\hline
Agence & AGENCE & référentiel des agences\\
\hline
Banque & B\_BANQUE & entités bancaires\\
\hline
BClient & B\_CLIENTS & clients enrichis (secteur, région)\\
\hline
BtkClient & CLIENT\_BTK & clients du réseau\\
\hline
Client & CLIENT & table client historique\\
\hline
Gestionnaire & GESTIONNAIRE & axe gestionnaire\\
\hline
Objectif & B\_OBJECTIF & production commerciale par période\\
\hline
Pointage & POINTAGE & présence : arrivée, départ, statut\\
\hline
DemandeAgence & DEMANDE\_AGENCE & demandes de création d'agence\\
\hline
DemandeEmploye & DEMANDE\_EMPLOYE & demandes de création d'employé\\
\hline
DemandeClient & DEMANDE\_CLIENT & demandes de création de client\\
\hline
DemandeModification & DEMANDES\_MODIFICATION & demandes de modification de profil\\
\hline
Notification & NOTIFICATION & messages destinés à un utilisateur\\
\hline
JournalAction & JOURNAL\_ACTION & journal d'audit\\
\hline
\end{tabularx}
